% TOP_PROBLEM: {"id": 30006745, "title": "The Generalized Sato-Tate Conjecture for Higher-Dimensional Abelian Varieties", "category_id": 1, "category": {"id": 1, "name": "number_theory", "display_name": "Number Theory", "description": "Properties of integers, prime numbers, Diophantine equations.", "slug": "number-theory", "order_index": 1, "created_at": "2026-07-31T15:26:25.670Z"}, "status": "open"}
% FIRST_POSTED: 2026-09-25
% ENTRY_KIND: statement_only
% !TeX program = lualatex
% Definition and sources only; this file is not an LLM solution attempt.
\documentclass[11pt]{article}
\usepackage[margin=25mm]{geometry}
\usepackage{fontspec}
\setmainfont{Latin Modern Roman}
\usepackage{amsmath,amssymb,mathtools}
\usepackage{unicode-math}
\setmathfont{Latin Modern Math}
\usepackage{xurl,hyperref}
\hypersetup{colorlinks=true,urlcolor=blue}
\setcounter{secnumdepth}{0}
\setlength{\parindent}{0pt}
\setlength{\parskip}{5pt}
\setlength{\emergencystretch}{3em}
\begin{document}
\section*{\texorpdfstring{The Generalized Sato-Tate Conjecture for Higher-Dimensional Abelian Varieties}{The Generalized Sato-Tate Conjecture for Higher-Dimensional Abelian Varieties}}\label{30006745}
\subsection{Definitions and mathematical statement}
For an abelian variety $A/K$ of dimension $g$ and a good finite place $v$, Frobenius on its first cohomology has eigenvalues of absolute value $(Nv)^{1/2}$. Divide them by this factor and use the resulting conjugacy class $c_v$ in the source's compact Sato--Tate group $\operatorname{ST}(A)$. The conjecture is
\[
 \frac1{\pi_K(x)}\sum_{Nv\le x,\ v\text{ good}}F(c_v)
 \longrightarrow\int_{\operatorname{ST}(A)}F(g)\,dg
\]
for every continuous class function $F$, with normalized Haar measure and $\pi_K(x)$ the number of prime ideals of norm at most $x$. The full group, including its components, and its Frobenius-class construction use the source convention.
\par\smallskip\textit{Formulation and concept references:} \href{https://doi.org/10.4007/annals.2011.173.3.12}{[S1]}.

\subsection{Short English statement}
Are normalized Frobenius symmetries of every abelian variety distributed according to Haar measure on its Sato--Tate group?
\subsection{Sources}
\begin{itemize}
\item[S1]T. Barnet-Lamb, D. Geraghty, M. Harris, R. Taylor, Ann. of Math. 173(3):1841-1869, 2011.
\url{https://doi.org/10.4007/annals.2011.173.3.12}.
\textit{Formulation source.}
\end{itemize}
\end{document}
